\section{An explicit bounded formula for every exterior-order index}
\label{sec:bounded-indices}

The tables of IF-03 can be extended without constructing a matrix on
the full exterior space and without an unbounded word-closure loop.
The following gives one finite integer matrix, of explicitly bounded
size, whose Smith factors are exactly all the desired order indices.
It is a formula for every rank, not a claim that a small-rank table
determines an unspecified theta period.

\begin{theorem}[SN-15: all-rank index matrices]
Let $V_n=\Lambda^*\mathbb Z^{2n}$, with its standard integral
symplectic basis, let $L$ be multiplication by its symplectic form,
and let $D$ be ordinary contraction as in IF-01. Put
\[
 R_n=\mathbb Q[L,D],\qquad
 A_n=R_n\cap\operatorname{End}_{\mathbb Z}(V_n),\qquad
 O_n=\mathbb Z[L,D],\qquad
 O_n^\pi=\mathbb Z[L,D,\pi_0,\ldots,\pi_{2n}],
\]
where $\pi_j$ is the exterior-degree projector. Let
\[
 \mathcal R_n=\{(v,k,l,t):0\le v\le n,\ 0\le k,l\le v,
 \ \max(0,k+l-v)\le t\le\min(k,l)\},
\]
and index columns by $(a,c,b)$ with $0\le a,c\le n$ and
$0\le b\le2n$. For $s=l-c$, define integer matrices $M_n$ and
$M_n^\pi$ by
\[
 (M_n)_{(v,k,l,t),(a,c,b)}=
 \begin{cases}
 a!c!\binom ts\,(v-2s)^b,&0\le s\le t,\quad k=s+a,\\
 0,&\text{otherwise},
 \end{cases}                                      \tag{SN-15a}
\]
and
\[
 (M_n^\pi)_{(v,k,l,t),(a,c,b)}=
 \begin{cases}
 a!c!\binom ts\,\mathbf1_{b=n-v+2l},
       &0\le s\le t,\quad k=s+a,\\
 0,&\text{otherwise}.
 \end{cases}                                      \tag{SN-15b}
\]
Here $0^0=1$. Each matrix has $\binom{n+4}{4}$ rows and
$(n+1)^2(2n+1)$ columns, and rank
$r_n=(n+1)(n+2)(2n+3)/6$.

The nonzero Smith factors of $M_n$ are precisely the invariant
factors of $A_n/O_n$ (including unit factors), and those of
$M_n^\pi$ are precisely the invariant factors of $A_n/O_n^\pi$.
In particular, if $\delta_j(M)$ is the nonnegative gcd of all
$j$-rowed minors, with $\delta_0=1$, then
\[
 [A_n:O_n]=\delta_{r_n}(M_n),\qquad
 [A_n:O_n^\pi]=\delta_{r_n}(M_n^\pi),\qquad
 d_j=\delta_j/\delta_{j-1}.                         \tag{SN-15c}
\]
These are explicit finite binomial-and-determinant formulas for
every rank. All their prime divisors are at most $n$.
\end{theorem}

\begin{proof}
For each symplectic pair choose one of its four exterior states
$1,e_i,f_i,e_i\wedge f_i$. Fix the pairs in singleton states
$e_i$ or $f_i$. On the remaining $v$ pairs, the basis is indexed
by subsets $S\subset\{1,\ldots,v\}$: a member of $S$ is in the
full state and a nonmember in the empty state. The actions are
\[
 L[S]=\sum_{i\notin S}[S\cup\{i\}],\qquad
 D[S]=\sum_{i\in S}[S\setminus\{i\}].             \tag{SN-15d}
\]
The signs are positive because the pair state has degree two;
the two contraction signs cancel. This is an integral direct-sum
decomposition, and every $v=0,\ldots,n$ occurs. Repeated copies
have identical operator matrices. We may therefore test matrix
integrality on one copy of each of these Boolean modules.

Set $H=DL-LD$. On $[S]$ it acts by $v-2|S|$; equivalently it
acts on exterior degree $j$ by $n-j$. Directly from (SN-15d),
\[
 [H,L]=-2L,\qquad [H,D]=2D,\qquad DL=LD+H.
\]
These integral commutation rules reorder every word in $L,D,H$
as an integral sum of $L^aH^bD^c$. For a detailed termination
argument, order the three letter types as $L,H,D$ and move each
out-of-order adjacent pair past one another. The exchanged term
has fewer inversions of that word length, while the commutator
terms have smaller word length. Induction on length and then
inversions terminates and introduces only integer coefficients.
Since $H\in O_n$, this gives a spanning set for $O_n$.
The identities $L^{n+1}=D^{n+1}=0$ and
\[
 \prod_{j=-n}^{n}(H-j)=0
\]
reduce it to $0\le a,c\le n$, $0\le b\le2n$; the polynomial
in $H$ is monic with integer coefficients. Thus the columns in
(SN-15a) will span the whole order, not just a finite selection
of its short words.

Fix a source subset $S$ of size $l$ and a target subset $T$ of
size $k$, with $|S\cap T|=t$. The intermediate subset after
$D^c$ has size $s=l-c$. For a nonzero contribution it must be a
subset of $S\cap T$, and $k=s+a$. There are $\binom ts$ such
intermediate subsets. Ordinary removal contributes $c!$, ordinary
addition contributes $a!$, and the intermediate action of $H^b$
contributes $(v-2s)^b$. This proves (SN-15a).

For the order with projectors, move every degree projector to
the right using $\pi_jL=L\pi_{j-2}$ and
$\pi_jD=D\pi_{j+2}$, with out-of-range projectors zero.
Orthogonality collapses their products. Reorder the remaining
word as above; $H^bD^c\pi_j=(n-j+2c)^bD^c\pi_j$.
Consequently $L^aD^c\pi_j$, $0\le a,c\le n$,
$0\le j\le2n$, span $O_n^\pi$ integrally. The source subset
has exterior degree $n-v+2l$, including the fixed singletons.
The preceding count proves (SN-15b).

All entries of a matrix from $R_n$ on a given Boolean module
depend only on $(|T|,|S|,|T\cap S|)$: permutations of the $v$
pairs commute with both operators and act transitively on pairs
of subsets having these three invariants. Every admissible
quadruple in $\mathcal R_n$ actually occurs. Recording one entry
per quadruple is therefore injective on $R_n$, and an element
of $R_n$ has an integral full exterior matrix if and only if
these recorded entries are integral. The count of quadruples
for a fixed $v$ is $\binom{v+3}{3}$: they are the four
nonnegative sizes of $S\cap T$, $T\setminus S$,
$S\setminus T$, and the complement. Summing gives the claimed
row count.

Thus in the recorded coordinate space, $A_n$ is exactly the
saturation of the column lattice of $M_n$. The same holds for
$M_n^\pi$: its rational span is still $R_n$, because IF-02
puts the degree projectors in $A_n\subset R_n$. IF-02 also
gives the stated rank. Smith normal form of an integer matrix
identifies the quotient of its saturated column span by its
column lattice with the direct sum of the cyclic groups of
its nonzero Smith factors. This proves all the index formulas.

Finally IF-02 gives
$A_n=\mathbb Z[L^{(a)},D^{(c)}:0\le a,c\le n]$.
After inverting $n!$, each generator is already in $O_n$,
so $A_n[1/n!]=O_n[1/n!]=(O_n^\pi)[1/n!]$.
The finite quotients can therefore have no prime divisor
larger than $n$. For $n=0$, all orders are $\mathbb Z$ and
the same formulas give index one.
\end{proof}

\paragraph{Integral row compression and independent shift blocks.}
For evaluation one may retain only the rows satisfying
\[
 t=0\quad\hbox{or}\quad v-k-l+t=0.                \tag{SN-15e}
\]
There are exactly $r_n$ of these rows. Removing the others preserves
all nonzero Smith factors, for the following integral reason.
Whenever $t>0$ and $v-k-l+t>0$, Pascal's identity in (SN-15a)
gives the row identity
\[
 R(v,k,l,t)=R(v,k,l,t-1)+R(v-2,k-1,l-1,t-1).      \tag{SN-15f}
\]
Both rows on the right are admissible. The intermediate powers of
$H$ agree, because replacing $(v,l)$ by $(v-2,l-1)$ leaves
$v-2(l-c)$ unchanged. The same argument works in (SN-15b), since
$n-(v-2)+2(l-1)=n-v+2l$. Recursive application strictly decreases
both positive quantities $t$ and $v-k-l+t$ along each nonterminal
branch. Every deleted row is therefore an integer combination of
retained rows, so subtracting these combinations gives zero rows
by unimodular row operations.

The four-size count in the proof shows that the number of retained
rows for a fixed $v$ is
$\binom{v+3}{3}-\binom{v+1}{3}$, where the subtracted term counts
the tuples with both the intersection and complement nonempty.
Summing gives
$\binom{n+4}{4}-\binom{n+2}{4}=r_n$. Since the original matrix
has rank $r_n$, the retained matrix has full row rank. Its
saturated column lattice is therefore the whole $\mathbb Z^{r_n}$.

Each nonzero entry further satisfies $k-l=a-c$. Sorting rows and
columns by this shift gives independent blocks for
$-n\le k-l\le n$, with no interactions between them. One may
compute each block's Smith factors, take the direct sum of the
resulting finite abelian groups, and sort the prime-wise valuations
to recover the global invariant factors. This is an integral block
decomposition, not a rational decomposition into irreducible
Lefschetz strings that might change the lattice.

\paragraph{Reproducible arithmetic and scope.}
The program \path{work/root/finite_index_formula.py} builds these
bounded matrices directly from (SN-15a)--(SN-15b) and records all
nonzero Smith factors. For $n=1,2,3,4$ it gives exactly the
indices of IF-03, for both orders; this is a third construction,
independent of the right-word and left-word closure programs.
It uses the same exact integer linear-algebra library and is
not described as an independent formal proof kernel.

For larger matrices, \path{work/root/finite_index_padic.py} performs
unit-pivot elimination modulo $p^c$ for each prime $p\le n$.
When no unit remains it divides the remaining matrix and modulus by
$p$; precision is doubled if exhausted, and no undecided exponent is
accepted. A unit pivot splits off one free summand over the local
ring; the division increases every remaining elementary-divisor
exponent by one. These two operations justify the computation.
The program \path{work/root/verify_finite_index_padic.py} checks the
new ranks instead by exact rational Schur complements, always
choosing a pivot of minimum $p$-adic valuation. All elimination
ratios are $p$-integral, so these are invertible row and column
operations over $\mathbb Z_p$. It uses neither modular precision
nor a Smith routine. The 144 block-and-prime computations at
$n=5,6$ agree, and the local method also reproduces every Smith
factor at $n=1,2,3,4$. The additional exact indices are
\[
\begin{array}{c|c|c|c}
n&r_n&[A_n:O_n]&[A_n:O_n^\pi]\\\hline
5&91&2^{201}3^{78}5^4&2^{80}3^{28}5^2\\
6&140&2^{420}3^{149}5^{20}&2^{168}3^{62}5^{10}.
\end{array}                                      \tag{SN-15g}
\]
All invariant factors, including unit factors, are recorded in the
JSON certificates. The prime-support proof above excludes any
unexamined prime; these are exact indices, not just lower bounds
from selected primes.

\paragraph{Geometric formulation.}
All matrices and formulas are defined over
$\operatorname{Spec}\mathbb Z$, before any rational or
$\ell$-adic realization. They concern the exterior model for
the Lefschetz algebra in cote 16 scans 014--018, with ordinary
contraction explicitly distinguished from normalized lowering.
Exterior powers give the $\lambda$-operations of SGA 6 V 2.2.1;
the Smith factors retain the integral information absent from
their rational $K_0$ character. Over $\mathbb Z_\ell$ the
quotients have lengths $v_\ell(\delta_{r_n})$. In the abelian
cohomological realization the rational Lefschetz operators are
algebraic by SN-10, but these arithmetic order indices alone
do not change a weight, determine a theta moment, or alter a
motivic local $L$-polynomial.
